\documentclass[11pt]{article}

\usepackage[utf8]{inputenc}
\usepackage[T1]{fontenc}
\usepackage{lmodern}
\usepackage[margin=2.5cm]{geometry}
\usepackage{amsmath,amssymb,amsthm}
\usepackage{hyperref}

\newtheorem{theorem}{Theorem}
\newtheorem{lemma}[theorem]{Lemma}
\newtheorem{corollary}[theorem]{Corollary}
\theoremstyle{definition}
\newtheorem{definition}[theorem]{Definition}

\title{Planar Graphs Are 6-Colorable}
\author{}
\date{}

\begin{document}

\maketitle

\section{Preliminaries}

All graphs are finite and simple (no loops, no parallel edges).
For a graph $G$ we write $V(G)$ and $E(G)$ for its vertex and edge sets and
$\deg(x)$ for the degree of a vertex $x$.

\begin{definition}
  A graph is \emph{planar} if it can be drawn in the plane such that edges
  are simple curves between their endpoints that meet only at common
  endpoints. Such a drawing is a \emph{plane graph}. The \emph{faces} of a
  plane graph are the connected components of the plane minus the drawing.
\end{definition}

\begin{definition}
  A \emph{$k$-coloring} of a graph $G$ is a map
  $c \colon V(G) \to \{1,\dots,k\}$ such that $c(x) \neq c(y)$ for every
  edge $xy \in E(G)$.
\end{definition}

Every subgraph of a planar graph is planar: simply delete the corresponding
points and curves from the drawing.

\section{Euler's Formula}

% lax begin Lax909950.EulerFormula
\begin{lemma}[Euler's formula]\label{lem:euler}
  Let $G$ be a connected plane graph with $v$ vertices, $e$ edges and
  $f$ faces. Then $v - e + f = 2$.
\end{lemma}
% lax end

% lax begin Lax909950Proofs.euler_formula
\begin{proof}
  We proceed by induction on $e$. Since $G$ is connected, $e \geq v - 1$.

  If $e = v - 1$, then $G$ is a tree. A drawing of a tree does not separate
  the plane, so $f = 1$ and $v - e + f = v - (v-1) + 1 = 2$.

  If $e \geq v$, then $G$ is not a tree and hence contains a cycle $C$.
  Let $uw$ be an edge of $C$. By the Jordan curve theorem, the drawing of $C$
  separates the plane into an inside and an outside region, and the two
  sides of the curve representing $uw$ lie in different faces of $G$.
  Deleting $uw$ therefore merges these two faces into one while leaving all
  other faces unchanged. The graph $G - uw$ is still connected (since $uw$
  lies on a cycle), and it has $v$ vertices, $e - 1$ edges and $f - 1$ faces.
  By induction, $v - (e-1) + (f-1) = 2$, that is, $v - e + f = 2$.
\end{proof}
% lax end

\section{Edge Density}

% lax begin Lax909950.EdgeDensity
\begin{lemma}[Edge density]\label{lem:density}
  Every planar graph $G$ with $v \geq 3$ vertices has at most $3v - 6$ edges.
\end{lemma}
% lax end

% lax begin Lax909950Proofs.edge_density
\begin{proof}
  Adding edges only increases $e$, so we may assume $G$ is edge-maximal
  planar; in particular $G$ is connected, since otherwise an edge joining
  two components could be added to the drawing without crossings. Fix a
  plane drawing of $G$ with $e$ edges and $f$ faces.

  Consider for each face its boundary walk and count the number of
  edge-sides it contains. Every edge has two sides, each of which lies on
  exactly one face boundary, so the total count is exactly $2e$. On the other
  hand, since $G$ is simple, connected and has $v \geq 3$ vertices (and thus
  $e \geq 2$), every face boundary contains at least $3$ edge-sides: a
  boundary with only one or two edge-sides would require a loop, a pair of
  parallel edges, or a graph consisting of a single edge. Hence
  $2e \geq 3f$, i.e.\ $f \leq \tfrac{2}{3}e$.

  By Euler's formula (Lemma~\ref{lem:euler}),
  \[
    2 = v - e + f \leq v - e + \tfrac{2}{3}e = v - \tfrac{1}{3}e,
  \]
  which rearranges to $e \leq 3v - 6$.
\end{proof}
% lax end

\section{A Vertex of Small Degree}

% lax begin Lax909950.LowDegree
\begin{lemma}[Low-degree vertex]\label{lem:lowdeg}
  Every nonempty planar graph $G$ has a vertex of degree at most $5$.
\end{lemma}
% lax end

% lax begin Lax909950Proofs.exists_degree_le_five
\begin{proof}
  Let $v = |V(G)|$ and $e = |E(G)|$. If $v \leq 6$, every vertex has degree
  at most $v - 1 \leq 5$. Otherwise $v \geq 3$, and by the edge density
  bound (Lemma~\ref{lem:density}),
  \[
    \sum_{x \in V(G)} \deg(x) = 2e \leq 6v - 12 < 6v.
  \]
  So the average degree is less than $6$, and some vertex has degree at
  most $5$. Equivalently: not all vertices can have high degree ($\geq 6$).
\end{proof}
% lax end

\section{The Six Color Theorem}

% lax begin Lax909950.SixColoring
\begin{theorem}[Six color theorem]\label{thm:six}
  Every planar graph admits a 6-coloring.
\end{theorem}
% lax end

% lax begin Lax909950Proofs.six_colorable
\begin{proof}
  By induction on the number of vertices $v$ of the planar graph $G$. If
  $v = 0$, the empty map is a 6-coloring.

  Let $v \geq 1$. By Lemma~\ref{lem:lowdeg}, $G$ has a vertex $x$ with
  $\deg(x) \leq 5$. The graph $G - x$ is a subgraph of $G$, hence planar,
  and has $v - 1$ vertices. By induction it has a 6-coloring $c$. The
  neighbors of $x$ use at most $5$ distinct colors under $c$, so some color
  $i \in \{1,\dots,6\}$ is not used by any neighbor of $x$. Extending $c$ by
  $c(x) = i$ yields a 6-coloring of $G$.
\end{proof}
% lax end

\end{document}
